% ============================================================
% VERSIÓN INDIVIDUAL — DANIEL DUARTE CORDERO
% TRACK A — datos tabulares.
% ============================================================
\section{Análisis Técnico del Track}
\label{sec:analisistecnico}

El proyecto trabaja en \textbf{Track A --- aprendizaje automático
clásico sobre datos tabulares}. \textbf{\% TODO integración:} insertar
la fecha exacta de aprobación comunicada por el docente.

\subsection{Análisis exploratorio}
\label{subsec:eda}
La auditoría preliminar contiene 52,085 episodios independientes y
7,038,378 snapshots frescos útiles. Se inspeccionaron 34 campos de
estado presentes en el 100\% de los snapshots muestreados. El target a
nivel de episodio se resuelve desde \texttt{rewards} en 99.86\% de los
casos: 27,135 episodios con recompensa $(1,-1)$, 24,844 con $(-1,1)$,
33 empates y 73 casos con final anómalo. La representación final desde
la perspectiva del jugador activo todavía debe materializarse para
medir su proporción de clases por snapshot; no se sustituye esa medida
por la distribución P0/P1 del episodio.

La muestra cubre diez días auditados y varias versiones del motor
(aproximadamente 1.30.1 a 1.32.4). Se identificaron 2,376 mazos y 3,139
agentes, con concentración apreciable y cambio temporal fuerte. La
unidad de agrupamiento es el episodio, no el snapshot. El EDA de la
tabla final reportará distribución por turno, episodios por día,
cardinalidad de categóricas, percentiles de numéricas y correlaciones
calculadas solo sobre train.

\subsection{Diccionario de datos}
\label{subsec:diccionario}
\begin{table}[!t]
\centering
\caption{Diccionario de datos (extracto).}
\label{tab:diccionario}
\footnotesize
\begin{tabularx}{\columnwidth}{@{}P{1.7cm} P{1.2cm} P{1.6cm} Y@{}}
\toprule
\textbf{Variable} & \textbf{Tipo} & \textbf{Origen} &
\textbf{Significado / rango / riesgo}\\
\midrule
\texttt{turn} & entero & \texttt{current.turn} & Progreso del juego; se usa para análisis por ventana y no como timestamp futuro.\\
\texttt{is\_first\_player} & binaria & \texttt{firstPlayer}/\texttt{yourIndex} & Orden de inicio; riesgo de atajo porque se confunde con asiento en días tardíos.\\
\texttt{prize\_remaining\_*} & entero & ranuras de \texttt{prize} & Premios restantes por jugador; feature derivada observable.\\
\texttt{active\_hp\_ratio} & continua & \texttt{active.hp/maxHp} & Salud relativa del activo; debe respetar $0\leq hp\leq maxHp$.\\
\texttt{energy\_diff} & entero & agregados self--opp & Diferencia de energías visibles en tablero.\\
\texttt{deck\_hash} & categórica & catálogo de mazos & Alto riesgo de memorizar metajuego; solo en ablación/contexto y con holdout por mazo.\\
\bottomrule
\end{tabularx}
\end{table}

El diccionario completo debe mantenerse en el repositorio e incluir
familias de progreso, recursos, tablero, energía, estados especiales,
diferenciales y contexto. Las columnas de agrupación y auditoría se
documentan, pero se separan del vector predictivo.

\subsection{Calidad de los datos}
\label{subsec:calidad}
La auditoría verificó 52,085 episodios válidos, cero errores de parseo,
una única firma de esquema top-level, \texttt{schema\_version=1}, cero
duplicados por \texttt{episode\_id}/hash canónico, cero solapamientos
entre días y cero self-play. Se identificaron cinco episodios declarados
en el manifest del 2026-08-04 pero ausentes de los archivos. Sobre la
tabla materializada se verificarán faltantes por columna, rangos
imposibles (por ejemplo, HP negativo o mayor a \texttt{maxHp}), conteos
negativos, \texttt{bench\_size>benchMax}, consistencia lógica y
cardinalidad de categóricas antes de cualquier modelado.

\subsection{Datos faltantes e imputación}
\label{subsec:imputacion}
Emmanuel et al. distinguen mecanismos de datos faltantes y muestran que
la eficacia de un método de imputación depende de las características
del dataset; en su experimento KNN y missForest intercambian cuál es
mejor según el conjunto de datos \cite{emmanuel_missing_2021}. Por ello,
la política del proyecto es medir y caracterizar la ausencia antes de
imputar, en vez de aplicar una técnica por defecto.

\noindent\textbf{Ausencia estructural.} En el replay existen valores
\texttt{None} que corresponden a estados semánticamente no aplicables u
ocultos por las reglas del juego, por ejemplo ausencia de un objeto o
información privada del rival. La auditoría observó \texttt{looking=None}
en 98.1\% de los estados inspeccionados. La decisión de codificar estos
casos semánticamente proviene del esquema y del dominio del juego;
\textbf{Paper 6 no formaliza por sí solo la categoría ``ausencia
estructural'' ni debe citarse como si lo hiciera}. Se requiere una
referencia adicional específica sobre missing-by-design/structural
missingness para sostener esa terminología en la versión final.

\noindent\textbf{Ausencia inesperada.} Si un campo requerido falta por
corrupción, cambio de esquema o error del parser, se medirá su
frecuencia. Si es residual, podrá excluirse la observación o episodio
con criterio documentado; si es sistemático, se tratará como defecto
del pipeline o cambio de esquema y se corregirá antes de entrenar, no
se ocultará mediante imputación.

\noindent\textbf{Estrategia adoptada.} Ninguna variable numérica se
imputa automáticamente. Primero se clasifica la causa de ausencia y el
tipo de variable; luego se decide si corresponde codificación
semántica, exclusión documentada o una técnica de imputación ajustada
solo con train. Esta postura es coherente con la conclusión de Emmanuel
et al. de que no existe un método universalmente superior
\cite{emmanuel_missing_2021}.

\subsection{Valores atípicos}
\label{subsec:atipicos}
La detección combina reglas de dominio y estadística descriptiva. Los
rangos imposibles se consideran errores potenciales; los extremos
válidos se auditan antes de decidir cualquier exclusión. Por ejemplo,
\texttt{n\_steps} tiene mediana 140 y máximo 6,900, mientras
\texttt{max\_turn} tiene mediana 13, P95 27 y máximo 1,035. Un episodio
de 1,035 turnos no se elimina por ser extremo: se investiga si
corresponde a timeout, bucle, interacción válida o error. Las reglas de
exclusión se fijan antes de observar el test.

\subsection{Desbalance de clases}
\label{subsec:desbalance}
A nivel de episodio, P0 gana 52.10\% y P1 47.70\%, con 0.063\% de
empates y 0.140\% de forfeits. Sin embargo, la etiqueta final se define
desde la perspectiva del jugador activo y cada bando aporta un número
distinto de snapshots; por eso la proporción de clases relevante debe
medirse sobre la tabla final antes de elegir una intervención.

Si el desbalance es leve, la opción inicial será no corregirlo. Van den
Goorbergh et al. muestran que, para regresión logística estándar y
ridge, RUS, ROS y SMOTE no mejoran sistemáticamente el AUROC y pueden
sobreestimar fuertemente la probabilidad de la clase minoritaria
\cite{vandenGoorbergh_imbalance_2022}. Dado que nuestra salida pretende
interpretarse como probabilidad, cualquier corrección debe evaluarse
también por calibración. Este artículo \textbf{no} avala preferir
\texttt{class\_weight}: esa técnica no fue evaluada y, por tanto, no se
adoptará por cita indirecta. Si el EDA revela un desbalance material,
se compararán alternativas únicamente dentro de train y se reportará
su efecto sobre discriminación y calibración.

\subsection{Prevención explícita de data leakage}
\label{subsec:leakage}
El riesgo principal del Track A en este proyecto proviene de la
multiplicidad de snapshots por episodio y de la presencia de información
que solo existe después de la decisión. La mitigación combina diseño de
particiones, exclusión de columnas y pruebas automáticas.

\subsubsection{Lista negra de información prohibida}
No pueden entrar al vector de características: \texttt{episode\_id};
nombres de equipo o agente; scores del manifest; timestamps como
predictores directos; \texttt{rewards} y resultado final; cualquier
campo posterior al snapshot; información privada del rival;
\texttt{visualize}; contenido futuro de logs; nombres de archivo, rutas
y hashes. En particular, \texttt{visualize} expone el mazo restante
ordenado y no representa información disponible para el agente en el
momento real de decisión.

\subsubsection{Taxonomía de riesgo}
\begin{table}[!t]
\centering
\caption{Vectores de fuga: riesgo, evidencia y mitigación.}
\label{tab:leakage}
\footnotesize
\begin{tabularx}{\columnwidth}{@{}P{1.9cm} Y P{2.2cm}@{}}
\toprule
\textbf{Tipo} & \textbf{Riesgo y evidencia medida} &
\textbf{Mitigación}\\
\midrule
Episode leakage & Un split aleatorio de snapshots coloca estados de la misma partida en train y test para prácticamente todos los episodios. & Split/CV por \texttt{episode\_id}; intersección de grupos debe ser vacía.\\
Temporal/target leakage & \texttt{rewards}, \texttt{visualize} o logs futuros revelan desenlace o información posterior. & Lista negra y parser limitado a \texttt{observation} hasta $t$.\\
Deck/agent leakage & IDs o metadata pueden actuar como proxies del metajuego/habilidad; top-10 de mazos concentra 39.94\%. & Metadata fuera de features; holdout por mazo/matchup y ablación del contexto de mazo.\\
\bottomrule
\end{tabularx}
\end{table}

\subsubsection{Verificación automática}
El pipeline incluirá al menos dos invariantes que deben fallar en CI o
antes de entrenar: (1) ninguna columna de la lista negra puede aparecer
en el conjunto de features, y (2) la intersección de
\texttt{episode\_id} entre cualquier par de particiones debe ser vacía.
Se añadirá una prueba del parser para confirmar que la información
privada del rival permanezca oculta y que ninguna transformación se
ajuste con datos de validación o test.
